\section{Synchronization}

\begin{example}

                Assume we have a program where two calculations are made at the same time using two diffrent threads. T1
                calculates "5+5" while T2 calculates "3+3" now what if we start both threads at once and they both write
                to the console first their calculation and then their answer. One possible interleaving would be "5+5
                3+3 = 9 = 10" or even worse "3+3 5+5 = 9 = 10"
            
\end{example}

\begin{example}

                Another famous example is the counter. When a counter c is incremented using c++ or ++c the value of c
                first
                gets read, then incremented and then stored. The counter fails if one thread stores his result between
                the time another reads and stores.
            
\end{example}

\begin{definition}

                Such a codeblock (incrementing the counter) that might lead to bad interleavings is called a critical
                section. To ensure program correctness we need the enforce mutual exclusion on them.
            
\end{definition}

\begin{definition}

                Thread safty refers to nothing bad happeing in all possible interleavings and thus implies program
                safety.
            
\end{definition}

\begin{definition}

                Thread Liveness refers to something good happening eventually, which implies the correctness of a
                program.
            
\end{definition}

\subsection{Locks and Synchronized}

            Every Java-Object has an intrinsic lock or lock monitor. These locks are unique and can be optaind by
            threads which allows us to lock a certain codeblock. If a thread wants to access such a codeblock and
            execute the code within it first has to optain the lock, while it has the lock, no other thread can accuire
            it.

            \begin{codeexample}

                synchronized(lockObject) {
                i++;
                }
            
\end{codeexample}

            This ensures mutual exlusion on the synchronized codeblock. It is also possible to use synchronized as a
            keyword in a function to sync. the whole function.

            \begin{remark}

                By default synchronized uses "this" as the lock-object.
            
\end{remark}

\begin{remark}

                Java-Locks are reentrant meaning they can be locked multiple times by the same Thread recursively (It
                will then have to release the look as many times as it aquired it).
            
\end{remark}

            For indementing and decrementing integers there is a special DataType AtomicInteger which has built in
            synchronization as incrementing/decrementing is done in a single step.

            \begin{codeexample}

                AtomicInteger counter = new AtomicInteger(0);
                counter.incrementAndGet();
            
\end{codeexample}

\subsection{Wait, notify, notifyAll}

            We can use wait() on a lock-object to tell the current thread to sleep until further notice if the lock is
            not available. This allows use to save ressources.

            \begin{codeexample}

                synchronized(lockObject) {
                while(!condition) {
                lockObject.wait();
                }
                // do work
                }
            
\end{codeexample}

            To wake one or multiple threads and tell them that the lock is now available we can use notify() or
            notifyAll() on the lock-object to wake up an aribtrary waiting thread or all of them.

            \begin{important}

                Always wait inside a while loop, as this ensures checking the condition after beeing notified.
            
\end{important}

\begin{codeexample}

                synchronized(lockObject) {
                while (!condition) {
                lockObject.wait();
                }
                //do work
                notifyAll();
                }
            
\end{codeexample}

\subsection{Lock API}

\begin{codeexample}

                class BankAccount {
                int balance() {return balance;}
                void setBalance (in x) {balance = x;}
                void withdraw(int amount) {
                int b = balance();
                if (b less than amount) {
                throw new Exception("Not enough money");
                }
                setBalance(b-amount);
                }
                }
            
\end{codeexample}

            The problem with the above code is that it is not thread-safe. If two threads try to withdraw money
            at the same time, they might both read the same balance and then both write to it, resulting in a
            wrong or even negative final balance. We need mutual exclusion in that critical section.
            \\
\\

            Instead of using built-in synchronized blocks we can also use locks from the java.util.concurrent.locks
            package.

            \begin{codeexample}

                class BankAccount {
                private Lock lock = new ReentrantLock();
                ...
                void withdraw(int amount) {
                lock.lock();
                try {
                ...
                } finally {
                lock.unlock();
                }
                }
                }
            
\end{codeexample}

\begin{important}

                Always use a try-finally block to ensure the lock is released even if an error occures.
            
\end{important}

\begin{remark}

                There is also a lock.tryLock(100, TimeUnit.MILLISECONDS) method which tries to accuire the lock only for
                a given amount of time. This is useful for avoiding deadlocks (ex. when two threads have each one lock
                and wait on the other to get released)
            
\end{remark}

\begin{tabular}{|c|c|c|}
\hline
Aspect & synchronized & Lock API \\\hline
Lock handling & Automatic (JVM managed) & Manual (lock() / unlock()) \\\hline
Lock release & Automatic & Must be done explicitly (finally!) \\\hline
Scope & Method or block scoped & Flexible (can span multiple methods) \\\hline
Waiting behavior & Thread blocks indefinitely & tryLock() avoids blocking \\\hline
Interruptibility & Not interruptible while waiting & lockInterruptibly() supports interrupt \\\hline
Flexibility & Low & High \\\hline
Complexity & Simple, less error-prone & More complex, error-prone \\\hline
Deadlock risk & Lower & Higher (if unlock forgotten) \\\hline
Typical use & Simple synchronization & Advanced concurrency control \\\hline
\end{tabular}

